\Askhsh{1707}{4}
\begin{ylh}{1}
    \item 3.6. Κριτήρια ισότητας ορθογώνιων τριγώνων.
\end{ylh}
Στο τρίγωνο $\mathrm{A}\mathrm{B}\Gamma$ ($\mathrm{A}\mathrm{B} < \mathrm{A}\Gamma$) του παρακάτω σχήματος, η κάθετη στο μέσο $\mathrm{M}$ της $\mathrm{B}\Gamma$ τέμνει την προέκταση της διχοτόμου $\mathrm{A}\Delta$ στο σημείο $\mathrm{E}$. Αν $\Theta, \mathrm{Z}$ είναι οι προβολές του $\mathrm{E}$ στις $\mathrm{A}\mathrm{B}, \mathrm{A}\Gamma$, να αποδείξετε ότι:
\begin{alist}
    \item Το τρίγωνο $\mathrm{E}\mathrm{B}\Gamma$ είναι ισοσκελές. \monades{5}
    \item Τα τρίγωνα $\Theta\mathrm{B}\mathrm{E}$ και $\mathrm{Z}\Gamma\mathrm{E}$ είναι ίσα. \monades{8}
    \item $\mathrm{A}\hat{\Gamma}\mathrm{E} + \mathrm{A}\hat{\mathrm{B}}\mathrm{E} = 180^\circ$. \monades{12}
\end{alist}
% TODO: Προσθήκη σχήματος
